
%%%%%%%%%%%%%%%%%%%%%%%%%%%%%%%%%%%%%%%%%%%%%%%%%%%%%%%%%%%%%
\section{模型假设}

为简化问题和突出研究重点，本文做出以下假设：

\begin{itemize}[itemindent=2em]
\item \textbf{假设1：数据完整性与真实性。} 假设输入系统的患者健康数据已经过医疗机构的初步验证，基本真实可信，经过数据清洗后能满足分析要求。对于缺失数据，采用合理的插补方法进行处理。

\item \textbf{假设2：患者特征的稳定性。} 假设在较短时间段内（如3个月以内），患者的基本特征（年龄、体质指数、既往史等）相对稳定，不考虑这些特征的长期演变。

\item \textbf{假设3：预测模型的有效期。} 假设基于历史数据训练的预测模型在应用期间仍然有效，当数据分布发生显著变化时需要重新训练（本文定期（每季度）进行模型更新）。

\item \textbf{假设4：患者群体的代表性。} 假设用于模型训练和验证的患者样本具有代表性，能够反映目标患者群体的特征分布。

\item \textbf{假设5：医患依从性。} 假设患者能够按照健康管理方案进行生活干预和定期复查，系统推荐的管理策略是科学合理的。

\item \textbf{假设6：隐私和安全。} 假设系统中患者信息的采集、存储、处理、使用均符合相关法律法规要求，建立了完善的隐私保护机制。

\item \textbf{假设7：技术环境的稳定性。} 假设系统运行所依赖的硬件、网络、数据库等基础设施保持稳定，不考虑突发故障情况。

\end{itemize}
